

% Para el titulo
\usepackage{titling}
\pretitle{\begin{center}
\href{https://www.ine.es}{\includegraphics[width=2in,height=2in]{fig/ine_logo.pdf}}\LARGE\\}
\posttitle{\end{center}}
% Para poder definir el background color de los chunks de R
% al generar el pdf
\usepackage{xcolor}
\definecolor{granate_ine}{RGB}{136,19,51}
\definecolor{verde_fuerte_ine}{RGB}{69,126,118}
\definecolor{verde_claro_ine}{RGB}{221,238,236}

\usepackage{etoolbox}
\patchcmd{\tableofcontents}{Contents}{\'Indice}{}{}

% Para los márgenes del pdf
\usepackage[a4paper,
            left=1in,
            right=1in,
            top=1.2in,
            bottom=1in,
            footskip=.25in]{geometry}
% Para cambiar el color de los títulos a granate INE
\usepackage{sectsty}
\chapterfont{\color{granate_ine}}  % sets colour of chapters
\sectionfont{\color{granate_ine}}  % sets colour of sections
\subsectionfont{\color{granate_ine}}  % sets colour of subsections


% Ponemos en español estas palabras
\renewcommand{\contentsname}{Indice}
\renewcommand{\chaptername}{Capítulo}

% Ponemos todo el texto en negrita como el verde INE
\let\oldtextbf\textbf
\renewcommand{\textbf}[1]{\textcolor{verde_fuerte_ine}{\oldtextbf{#1}}}


